% TOP_PROBLEM: {"problemId": "problem.coleman-oort-conjecture-on-shimura-subvarieties-of-the-torelli-locus", "canonicalTitle": "Coleman-Oort conjecture: eventual special-subvariety exclusion", "releaseRank": 403, "primaryDomain": "number_theory_arithmetic_geometry", "primaryDomainLabel": "Number theory & arithmetic geometry", "displayStatus": "open_with_solved_subcases", "releaseStatus": "open_with_solved_subcases"}
% !TeX program = lualatex
% ATTEMPT_STATUS: unresolved
% Research continuation by GPT_6_Astra_Ultra, 27 September 2026.
\documentclass[11pt]{article}
\usepackage[margin=25mm]{geometry}
\usepackage{fontspec}
\setmainfont{Latin Modern Roman}
\usepackage{amsmath,amssymb,mathtools}
\usepackage{unicode-math}
\setmathfont{Latin Modern Math}
\usepackage{xurl,hyperref}
\hypersetup{colorlinks=true,urlcolor=blue}
\setcounter{secnumdepth}{0}
\setlength{\parindent}{0pt}
\setlength{\parskip}{5pt}
\setlength{\emergencystretch}{3em}
\begin{document}
\section*{\texorpdfstring{Coleman-Oort conjecture: eventual special-subvariety exclusion}{Coleman-Oort conjecture: eventual special-subvariety exclusion}}\label{403}
\subsection{Definitions and mathematical statement}
Let $\mathcal A_g$ be the complex moduli space of principally polarized abelian varieties of dimension $g$. The Torelli map sends a smooth genus-$g$ curve to its Jacobian with its principal polarization; denote its image by $\mathcal T_g^\circ$ and its Zariski closure by $\mathcal T_g$. A special subvariety is an irreducible Shimura subvariety in the usual Siegel-moduli sense. The Coleman--Oort assertion is
\[
 \exists g_0\ \forall g\ge g_0:\quad
 \nexists S\subseteq\mathcal A_g\text{ special with }
 \dim S>0,\ S\subseteq\mathcal T_g,\ S\cap\mathcal T_g^\circ\ne\varnothing.
\]
Thus the prohibited subvariety is generically made of smooth Jacobians. No value or effectiveness of $g_0$ is requested. Zero-dimensional CM points and special subvarieties wholly in the decomposable boundary are not excluded by this statement.
\par\smallskip\textit{Formulation and concept references:} \href{https://arxiv.org/pdf/2102.12349v2}{[S1]}.

\subsection{Short English statement}
In sufficiently large genus, are positive-dimensional special subvarieties excluded from the Torelli locus whenever they contain smooth Jacobians generically?
\subsection{Research attempt}
\textit{GPT-6 Astra Ultra; 27 September 2026. Status: UNRESOLVED.}

I test the two most accessible approaches: dimension comparison and
families of covers. The calculations show a growing numerical
obstruction in a concrete cyclic-cover family. I also construct
positive-dimensional special subvarieties in the decomposable Torelli
boundary in every genus $g\ge2$, verifying exactly why they are not
counterexamples. The smooth-Jacobian intersection condition is essential
throughout.

\paragraph{Codimension alone is insufficient.}
For $g\ge2$, the dimensions of curve moduli and Siegel moduli are
$3g-3$ and $g(g+1)/2$, respectively. Torelli's theorem gives the
Torelli locus the former dimension. Hence its codimension is
\[
 \frac{g(g+1)}2-(3g-3)=\frac{(g-2)(g-3)}2.            \tag{1}
\]
This grows quadratically, making containment of a large special
subvariety restrictive. But the conjecture excludes even special
curves, not just subvarieties whose dimension is close to the ambient
dimension. A subvariety of arbitrarily large codimension can contain
a line: the coordinate subspace
$\mathbb A^r\times\{0\}\subset\mathbb A^{r+c}$ contains coordinate
lines for every $c$. Thus the growth in (1) cannot by itself prove
the required emptiness. One needs information about which tangent
directions or Hodge structures can lie in the Torelli locus.

\paragraph{A boundary stress test for every $g\ge2$.}
Fix a CM elliptic curve $E_0$ and vary another elliptic curve $E$.
The product principally polarized abelian varieties
\[
 E\times E_0^{g-1}                                    \tag{2}
\]
trace a positive-dimensional special subvariety of $\mathcal A_g$:
it is the image of the elliptic Shimura curve times fixed CM points
under the product morphism of Siegel moduli. One may work with a
level structure to make this morphism unambiguous, then take its
irreducible image after forgetting the level. The image has dimension
one, since varying $E$ gives infinitely many isomorphism classes.

These points lie in $\mathcal T_g$. Glue $g$ elliptic curves, the first
equal to $E$ and the rest to $E_0$, in a chain at marked points.
The resulting stable curve has compact type: its dual graph is a
tree. The normalization sequence for line bundles shows that its
connected Picard variety is the product of those of its components;
the possible torus factor has rank equal to the first Betti number
of the dual graph, which is zero. Its principally polarized Jacobian
is therefore (2). Smoothing the nodes gives smooth curves of genus
$g$, so these product Jacobians are limits of smooth Jacobians in
$\mathcal A_g$. These standard compact-type Torelli facts and the
special-subvariety interpretation are recorded in [E1].

For $g\ge2$, none of (2) is the principally polarized Jacobian of a
smooth curve. The theta divisor of a smooth Jacobian is the image of
the irreducible variety $\operatorname{Sym}^{g-1}C$, up to translation,
and is therefore irreducible. The product polarization of two
positive-dimensional factors has theta divisor
\[
 (\Theta_1\times A_2)\ \cup\ (A_1\times\Theta_2),
\]
which is reducible. This property is preserved by polarized
isomorphism. Consequently the special curve constructed from (2) is
wholly in the decomposable boundary and misses
$\mathcal T_g^\circ$. It fails an explicit hypothesis of the proposed
counterexample, rather than contradicting the conjecture.

\paragraph{A family with computable Hodge dimensions.}
Consider cyclic triple covers of the projective line
\[
 C:\quad y^3=\prod_{i=1}^{3k}(x-t_i),
 \qquad t_i\ne t_j\ (i\ne j),\quad k\ge2.            \tag{3}
\]
Take the smooth projective normalization. Each finite branch point
has ramification index three, and the three points over infinity are
unramified because the polynomial degree is divisible by three.
Riemann--Hurwitz gives
\[
 2g-2=3(-2)+(3k)(3-1),\qquad g=3k-2.                \tag{4}
\]
The unordered branch configuration has dimension $3k-3$ after
quotienting by projective automorphisms. For $g\ge2$ the map to curve
moduli has finite fibers: a fixed curve has finite automorphism group,
so only finitely many possible order-three deck subgroups and quotient
maps up to the relevant equivalences. Torelli then gives the same
dimension to the Jacobian family.

The two nontrivial eigenspaces for $y\mapsto\zeta_3 y$ can be
calculated directly. For $j=1,2$, the forms
\[
 x^a\frac{dx}{y^j},\qquad 0\le a\le kj-2,
                                                        \tag{5}
\]
are holomorphic. Near a finite branch point, use $x-t_i=u^3$ and
$y$ equal to $u$ times a unit; the differential has order $2-j\ge0$.
At infinity, $y$ has pole order $k$ and $dx$ pole order two, so the
order is $kj-a-2\ge0$. The forms in each family are linearly
independent, and the two characters are distinct. Their total number
is $(k-1)+(2k-1)=3k-2=g$, proving they give the entire holomorphic
differential space. The eigenspace dimensions are therefore
\[
 d_1=k-1,\qquad d_2=2k-1.                            \tag{6}
\]

The polarized Hodge structures with this imaginary quadratic action
lie in the corresponding unitary PEL locus. Its symmetric domain has
dimension $d_1d_2$: infinitesimally the moving positive subspace has
tangent space $\operatorname{Hom}(\mathbb C^{d_1},\mathbb C^{d_2})$.
The difference between this dimension and the family dimension is
\[
 (k-1)(2k-1)-(3k-3)=2(k-1)(k-2).                    \tag{7}
\]
For $k=2$ the dimensions agree; for every $k\ge3$ the cover family
has strictly smaller dimension than this ambient PEL locus. The
comparison is exact and tends in the expected direction, but its
logical scope must be respected.

\paragraph{What the dimension obstruction actually proves.}
If the smallest special subvariety containing the family (3) is the
full PEL component just described, then (7) proves that the family
itself is not special for $k\ge3$. Indeed a special family would be
its own special closure, of the smaller dimension. Establishing that
closure, for example by determining the generic Mumford--Tate group,
is additional input not proved here. Merely noticing an action of
$\mathbb Q(\zeta_3)$ does not exclude further Hodge tensors that
could shrink the special closure.

Even a proof that the entire cover family is not special would not
exclude a smaller special curve lying inside it. And excluding such
curves in this one family would not control arbitrary Jacobians in
the Torelli locus. These are two separate quantifier gaps. Thus the
calculation does not extend a classification of selected covering
families to the universal assertion in the question.

\paragraph{Sources, verification, and outcome.}
[S1] studies evidence through Galois-cover families and their numerical
conditions, with explicit restrictions on the families covered. The
paper is not a classification of all possible positive-dimensional
special subvarieties of every Torelli locus. The finite arithmetic
check here verifies (4), the form counts (6), and the dimension
difference (7) for a range of $k$; the displayed formulas prove those
identities for all $k$. No monodromy computation or classification is
claimed.

The strongest results of the attempt are the exact boundary
countertest and the cyclic-cover Hodge calculation. The first gap in
the latter route is to determine the special closure, followed by the
much stronger task of excluding every smaller special subvariety
meeting the smooth locus in all sufficiently large genera. No
uniform genus threshold or general obstruction has been obtained.
The Coleman--Oort assertion remains \textbf{UNRESOLVED}.

\subsection{Sources}
\begin{itemize}
\item[S1]Diego Conti, Alessandro Ghigi and Roberto Pignatelli, Some evidence for the Coleman-Oort conjecture, arXiv2102.12349v2.
\url{https://arxiv.org/pdf/2102.12349v2}.
\textit{Formulation source.}
\item[E1]Ben Moonen and Frans Oort, \emph{The Torelli locus and special
subvarieties}, Handbook of Moduli II; author manuscript.
\url{https://www.math.ru.nl/personal/bmoonen/Papers/SpecialSubvarsTorelli.pdf}.
\textit{Primary author exposition for the compact-type boundary,
Jacobian indecomposability, and special loci with extra endomorphisms.
The cyclic-cover form counts are derived explicitly in this attempt.
Consulted 27 September 2026.}
\item[E2]Diego Conti, Alessandro Ghigi, and Roberto Pignatelli,
\emph{Some evidence for the Coleman--Oort conjecture}, arXiv:2102.12349v2.
\url{https://arxiv.org/pdf/2102.12349v2}.
\textit{Scope check of inherited [S1]: numerical conditions on specified
covering families are distinguished from the full conjecture.
Consulted 27 September 2026.}

\end{itemize}
\end{document}
